% !TEX root = preprint.tex
% Generated from manuscript.docx; see data/manuscript_docx/README.md.
\newcommand{\papertitle}{Near real-time data on the human neutralizing antibody landscape to influenza virus in summer of 2026 shows antigenic advance of H3N2 subclade K region D mutants and H1N1 D.3.1.1 Sa mutants}
\newcommand{\paperauthors}{Caroline Kikawa\textsuperscript{1,2,3,†}, Andrew Butler\textsuperscript{1,†}, John Huddleston\textsuperscript{4}, Sam A. Turner\textsuperscript{5}, Heidi Peck\textsuperscript{6}, Janet A. Englund\textsuperscript{7}, Kirsten Lacombe\textsuperscript{7}, Michael Busch\textsuperscript{8,9}, Marion C. Lanteri\textsuperscript{9,10}, Mars Stone\textsuperscript{8,9}, Bryan Spencer\textsuperscript{11}, Alexander L. Greninger\textsuperscript{12}, Derek J. Smith\textsuperscript{5}, Stephanie Wallace\textsuperscript{13}, Helen S. Marshall\textsuperscript{14}, Shidan Tosif\textsuperscript{15}, Scott E. Hensley\textsuperscript{16}, Ian G. Barr\textsuperscript{6}, Jesse D. Bloom\textsuperscript{1,2,17,\#,*}}
\newcommand{\paperaffiliations}{%
\textsuperscript{1}Division of Basic Sciences and Computational Biology Program, Fred Hutch Cancer Center, 1100 Fairview Ave N, Seattle, WA 98109, USA\\
\textsuperscript{2}Department of Genome Sciences, University of Washington, 1959 NE Pacific St, Seattle, WA 98195, USA\\
\textsuperscript{3}Medical Scientist Training Program, University of Washington, 1959 NE Pacific St, Seattle, WA 98195, USA\\
\textsuperscript{4}Vaccine and Infectious Disease Division, Fred Hutch Cancer Center, 1100 Fairview Ave N, Seattle, WA 98109, USA\\
\textsuperscript{5}Center for Pathogen Evolution, Department of Zoology, University of Cambridge, The Old Schools, Trinity Ln, Cambridge CB2 1TN, United Kingdom\\
\textsuperscript{6}WHO Collaborating Centre for Reference and Research on Influenza, The Peter Doherty Institute for Infection and Immunity, 792 Elizabeth Street Melbourne, 3000, Australia\\
\textsuperscript{7}Seattle Children's Research Institute and Department of Pediatrics, University of Washington, 1959 NE Pacific St, Seattle, WA 98195, USA\\
\textsuperscript{8}Vitalant Research Institute, 360 Spear St Ste 200, San Francisco, CA 94105, USA\\
\textsuperscript{9}Department of Laboratory Medicine, University of California San Francisco, 505 Parnassus Avenue, San Francisco, CA 94143, USA\\
\textsuperscript{10}Creative Testing Solutions, 2424 W Erie Dr, Tempe, AZ 85282, USA\\
\textsuperscript{11}American Red Cross, 180 Rustcraft Rd, Dedham, MA 02026, USA\\
\textsuperscript{12}Department of Laboratory Medicine and Pathology, University of Washington Medical Center, 1959 NE Pacific St, Seattle, WA 98195, USA\\
\textsuperscript{13}University of the Sunshine Coast Clinical Trials, Sunshine Coast, Queensland, Australia\\
\textsuperscript{14}Adelaide University, South Australia, 5005, Australia\\
\textsuperscript{15}Murdoch Children's Research Institute, University of Melbourne, The Royal Children's Hospital Melbourne, 50 Flemington Road, Parkville, Victoria, 3052, Australia\\
\textsuperscript{16}Depts of Microbiology and Medicine, Perelman School of Medicine, University of Pennsylvania, 3451 Walnut Street, Philadelphia, PA 19104, USA\\
\textsuperscript{17}Howard Hughes Medical Institute, 1100 Fairview Ave N, Seattle, WA 98109, USA\\
\textsuperscript{†} Authors contributed equally\\
\textsuperscript{\#} Lead contact\\
* Correspondence: \href{mailto:jbloom@fredhutch.org}{jbloom@fredhutch.org}
}
\newcommand{\paperabstract}{%
Human seasonal influenza evolves rapidly, necessitating twice yearly decisions about whether to update the strains in the vaccine. To help inform this decision, we have been using high-throughput sequencing-based neutralization assays to make twice yearly measurements of how recent human sera neutralize current human H3N2 and H1N1 strains. Here we provide the third installment in this series of measurements by reporting 47,851 titers representing neutralization of 148 viral strains by 325 human sera collected between April and August of 2026. Our measurements show that new H3N2 subclade K strains with mutations in antigenic region D and new H1N1 subclade D.3.1.1 strains with mutations in antigenic region Sa (such as G155E) have reduced neutralization by human sera, with notable heterogeneity in the impact of some of these mutations across sera from different individuals. This paper is accompanied by an interactive summary (\href{https://jbloomlab.github.io/flu-seqneut-2026/summary.html}{https://jbloomlab.github.io/flu-seqneut-2026/summary.html}) that enables detailed exploration of the results, and all titer data are publicly available for further analysis to aid vaccine antigen selection and studies of viral evolution.
}
\newcommand{\papershortauthors}{Kikawa, Butler et al.}
\newcommand{\papershorttitle}{Human neutralizing antibody landscape to influenza, summer 2026}
